\documentclass[multi=tikzpicture,crop]{standalone}

\input{../typography.tex.preamble}
\input{../tikz.tex.preamble}
\renewcommand{\arcsin}{\sin^{-1}}
\renewcommand{\arccos}{\cos^{-1}}
\renewcommand{\arctan}{\tan^{-1}}

\pgfplotsset{
  trigs/.style = {
    basic curve,
    xmin={-1.5*pi-pi/6}, xmax={1.5*pi+pi/6},
    xlabel=\empty, ylabel=\empty,
    xtick={-2*pi, -3*pi/2, -pi, -pi/2, 0, pi/2, pi, 3*pi/2, 2*pi},
    xticklabels={\(-2\pi\),,\(-\pi\),,,,\(\pi\),,\(2\pi\)},
    grid=none,
    samples=500,
    domain={-2*pi-pi/6}:{2*pi+pi/6},
  },
  sincos/.style = {
    trigs,
    ymin={-1.2}, ymax={1.2},
  },
  tansec/.style = {
    trigs,
    ymin={-1.2}, ymax={1.2},
    restrict y to domain=-5.25:5.25,
    ymin={-3.25}, ymax={3.25},
    ytick={-5, -4, -3, -2,-1,0,1, 2, 3, 4, 5},
    yticklabels={,-4,,-2,,0,,2,,4,},
  }
}

\pgfplotsset{
  inverse trigs/.style = {
    basic curve,
    xlabel={\(x\)}, ylabel={\(\theta\)},
    xlabel style = {anchor=west},
    ylabel style = {anchor=west},
    xtick distance={2},
    grid=none,
    samples=500,
  },
  arcsin/.style = {
    inverse trigs,
    xmin={-1.2}, xmax={1.2},
    xtick={-1,0,1},
    ytick={-pi, -pi/2, 0, pi/2, pi},
    yticklabels={\(-\pi\), \(-\pi/2\), 0, \(\pi/2\), \(\pi\)},
    domain={-1}:{1},
  },
  arctan/.style = {
    inverse trigs,
    xmin={-3.5}, xmax={3.5},
    ytick={-pi, -pi/2, 0, pi/2, pi},
    yticklabels={\(-\pi\), \(-\pi/2\), 0, \(\pi/2\), \(\pi\)},
    domain={-3.5}:{3.5},
  },
}

% Every plot comes in two versions: one with a title naming the function and
% one without.  Use an untitled page whenever the name of the function is the
% answer to the exercise (see lessons/review-trigs.tex).
%
% The inverse trig curves (arcsin, arccos, arctan, arccot) are drawn ultra
% thick in black, so black reads as "the graph" and supp purple as "the range".
% Every arcsin and arccos page (3, 6, 15, 16, 29, 30) also carries solid dots
% at x = -1 and x = 1, since those endpoints belong to the graph.
%
% The titled inverse trig plots (3, 6, 9, 12 and the aligned trio 29--31) also
% carry their range on the vertical axis, drawn in supp with filled/hollow
% endpoints exactly like the number lines in the notes, so the graph alone
% states the domain (in the title) and the range.  The untitled copies
% (15, 16, 17, 26) stay plain: those are the ones students annotate.
%
% PAGE MAP -- pages 1--17 are unchanged.  Other files index into this file by
% page number (slides.tex and lessons/review-trigs.tex), so append new plots at
% the end rather than inserting them.
%
%       plot                     titled   untitled
%       sin                        1        18
%       sin restricted             2        19
%       arcsin                     3        15
%       cos                        4        20
%       cos restricted             5        21
%       arccos                     6        16
%       tan                        7        22
%       tan restricted             8        23
%       arctan                     9        17
%       cot                       10        24
%       cot restricted            11        25
%       arccot                    12        26
%       sec                       27        13
%       csc                       28        14
%
% Pages 29--31 are arcsin, arccos and arctan again, drawn on one common
% y-range so that the three sit side by side with their x-axes at the same
% height (see the reference trio in lessons/review-trigs.tex).  Keep their
% ymin/ymax identical to each other, or the alignment is lost.

% The range of an inverse trig function, drawn on the vertical axis: a
% supp-coloured segment, with a filled dot at an endpoint that belongs to the
% range and a hollow dot at one that does not.
%
%   \rangeseg{<from>}{<to>}   \rangedot{<y>}   \rangehole{<y>}
\newcommand{\rangeseg}[2]{%
  \addplot[supp, line width=1.4pt, no markers] coordinates {(0,#1) (0,#2)};}
\newcommand{\rangedot}[1]{%
  \addplot[supp, only marks, mark=*, mark size=2.2pt] coordinates {(0,#1)};}
% Solid dots at the ends of a curve, marking endpoints that belong to the
% graph (arcsin and arccos at x = -1 and x = 1).
%
%   \curveends{<x1>}{<y1>}{<x2>}{<y2>}
\newcommand{\curveends}[4]{%
  \addplot[black, only marks, mark=*, mark size=2.2pt]
    coordinates {(#1,#2) (#3,#4)};}
\newcommand{\rangehole}[1]{%
  \addplot[supp, only marks, mark=*, mark size=2.2pt,
           mark options={fill=white, draw=supp, line width=0.9pt}]
    coordinates {(0,#1)};}

\begin{document}
%%% ------------------------------------------------------------ pages 1--12
%%% page 1: sin
\begin{tikzpicture}
  \begin{axis}[sincos, title={\(\sin(x)\)}]
    \addplot[very thick] {sin(x)};
  \end{axis}
\end{tikzpicture}

%%% page 2: sin restricted
\begin{tikzpicture}
  \begin{axis}[
    sincos,
    title={\(\sin(x)\) restricted to domain \([-\pi/2,\pi/2]\)},
    xtick={-pi/2, pi/2},
    xticklabels={\(-\pi/2\), \(\pi/2\)},
    ]
    \addplot[very thick, main, domain={-pi/2}:{pi/2}] {sin(x)};
  \end{axis}
\end{tikzpicture}

%%% page 3: arcsin
\begin{tikzpicture}
  \begin{axis}[
    arcsin,
    title={\(\theta = \arcsin(x)\) on domain \([-1,1]\)},
    ymin={-pi/2-pi/6}, ymax={pi/2+pi/6},
    ]
    \addplot[ultra thick, black] {asin(x)};
    \rangeseg{-pi/2}{pi/2}
    \rangedot{-pi/2}
    \rangedot{pi/2}
    \curveends{-1}{-pi/2}{1}{pi/2}
  \end{axis}
\end{tikzpicture}

%%% page 4: cos
\begin{tikzpicture}
  \begin{axis}[sincos, title={\(\cos(x)\)}]
    \addplot[very thick] {cos(x)};
  \end{axis}
\end{tikzpicture}

%%% page 5: cos restricted
\begin{tikzpicture}
  \begin{axis}[sincos, title={\(\cos(x)\) restricted to domain \([0,\pi]\)}]
    \addplot[very thick, main, domain={0}:{pi}] {cos(x)};
  \end{axis}
\end{tikzpicture}

%%% page 6: arccos
\begin{tikzpicture}
  \begin{axis}[
    arcsin,
    title={\(\theta = \arccos(x)\) on domain \([-1,1]\)},
    ymin={0-pi/6}, ymax={pi+pi/6},
    ]
    \addplot[ultra thick, black] {acos(x)};
    \rangeseg{0}{pi}
    \rangedot{0}
    \rangedot{pi}
    \curveends{-1}{pi}{1}{0}
  \end{axis}
\end{tikzpicture}

%%% page 7: tan
\begin{tikzpicture}
  \begin{axis}[tansec, title={\(\tan(x)\)}]
    \addplot[very thick] {tan(x)};
    \addplot[dotted] coordinates {(-1*pi/2, 5)  (-1*pi/2, -5)};
    \addplot[dotted] coordinates {(-3*pi/2, 5)  (-3*pi/2, -5)};
    \addplot[dotted] coordinates {( 1*pi/2, 5)  ( 1*pi/2, -5)};
    \addplot[dotted] coordinates {( 3*pi/2, 5)  ( 3*pi/2, -5)};
  \end{axis}
\end{tikzpicture}

%%% page 8: tan restricted
\begin{tikzpicture}
  \begin{axis}[
    tansec,
    title={\(\tan(x)\) with domain restriction \((-\pi/2,\pi/2)\)},
    xtick={-pi/2, pi/2},
    xticklabels={\(-\pi/2\), \(\pi/2\)},
    ]
    \addplot[very thick, main, domain={-pi/2}:{pi/2}] {tan(x)};
    \addplot[dotted, supp] coordinates {(-1*pi/2, 5)  (-1*pi/2, -5)};
    \addplot[dotted, supp] coordinates {( 1*pi/2, 5)  ( 1*pi/2, -5)};
  \end{axis}
\end{tikzpicture}

%%% page 9: arctan
\begin{tikzpicture}
  \begin{axis}[
    arctan,
    title={\(\theta = \arctan(x)\) on domain \((-\infty,\infty)\)},
    ymin={-pi/2-pi/6}, ymax={pi/2+pi/6},
    ]
    \addplot[ultra thick, black] {atan(x)};
    \addplot[dotted, supp] coordinates {(-10, pi/2)  ( 10, pi/2)};
    \addplot[dotted, supp] coordinates {(-10,-pi/2)  ( 10,-pi/2)};
    \rangeseg{-pi/2}{pi/2}
    \rangehole{-pi/2}
    \rangehole{pi/2}
  \end{axis}
\end{tikzpicture}

%%% page 10: cot
\begin{tikzpicture}
  \begin{axis}[tansec, title={\(\cot(x)\)}]
    \addplot[very thick] {cot(x)};
    \addplot[dotted] coordinates {(-2*pi, 5)  (-2*pi, -5)};
    \addplot[dotted] coordinates {(-1*pi, 5)  (-1*pi, -5)};
    \addplot[dotted] coordinates {( 1*pi, 5)  ( 1*pi, -5)};
    \addplot[dotted] coordinates {( 2*pi, 5)  ( 2*pi, -5)};
  \end{axis}
\end{tikzpicture}

%%% page 11: cot restricted
\begin{tikzpicture}
  \begin{axis}[
    tansec,
    title={\(\cot(x)\) restricted to domain \((0,\pi)\)},
    xtick={0,pi},
    xticklabels={\(0\), \(\pi\)}
    ]
    \addplot[very thick, main, domain={0}:{pi}] {cot(x)};
    \addplot[dotted, supp] coordinates {( 1*pi, 5)  ( 1*pi, -5)};
  \end{axis}
\end{tikzpicture}

%%% page 12: arccot
\begin{tikzpicture}
  \begin{axis}[
    arctan,
    title={\(\theta = \cot^{-1}(x)\) on domain \((-\infty,\infty)\)},
    ymin={0-pi/6}, ymax={pi+pi/6},
    ]
    \addplot[ultra thick, black] {pi/2-atan(x)};
    \addplot[dotted, supp] coordinates {(-10, 0)  ( 10, 0)};
    \addplot[dotted, supp] coordinates {(-10,pi)  ( 10,pi)};
    \rangeseg{0}{pi}
    \rangehole{0}
    \rangehole{pi}
  \end{axis}
\end{tikzpicture}

%%% ----------------------------------------------------------- pages 13--17
%%% page 13: sec, untitled
\begin{tikzpicture}
  \begin{axis}[tansec]
    \addplot[very thick] {1/cos(x)};
    \addplot[dotted] coordinates {(-1*pi/2, 5)  (-1*pi/2, -5)};
    \addplot[dotted] coordinates {(-3*pi/2, 5)  (-3*pi/2, -5)};
    \addplot[dotted] coordinates {( 1*pi/2, 5)  ( 1*pi/2, -5)};
    \addplot[dotted] coordinates {( 3*pi/2, 5)  ( 3*pi/2, -5)};
  \end{axis}
\end{tikzpicture}

%%% page 14: csc, untitled
\begin{tikzpicture}
  \begin{axis}[tansec]
    \addplot[very thick] {1/sin(x)};
    \addplot[dotted] coordinates {(-1*pi, 5)  (-1*pi, -5)};
    \addplot[dotted] coordinates {(-2*pi, 5)  (-2*pi, -5)};
    \addplot[dotted] coordinates {( 1*pi, 5)  ( 1*pi, -5)};
    \addplot[dotted] coordinates {( 2*pi, 5)  ( 2*pi, -5)};
  \end{axis}
\end{tikzpicture}

%%% page 15: arcsin, untitled
\begin{tikzpicture}
  \begin{axis}[
    arcsin,
    ymin={-pi/2-pi/6}, ymax={pi/2+pi/6},
    ]
    \addplot[ultra thick, black] {asin(x)};
    \curveends{-1}{-pi/2}{1}{pi/2}
  \end{axis}
\end{tikzpicture}

%%% page 16: arccos, untitled
\begin{tikzpicture}
  \begin{axis}[
    arcsin,
    ymin={0-pi/6}, ymax={pi+pi/6},
    ]
    \addplot[ultra thick, black] {acos(x)};
    \curveends{-1}{pi}{1}{0}
  \end{axis}
\end{tikzpicture}

%%% page 17: arctan, untitled
\begin{tikzpicture}
  \begin{axis}[
    arctan,
    ymin={-pi/2-pi/6}, ymax={pi/2+pi/6},
    ]
    \addplot[ultra thick, black] {atan(x)};
    \addplot[dotted, supp] coordinates {(-10, pi/2)  ( 10, pi/2)};
    \addplot[dotted, supp] coordinates {(-10,-pi/2)  ( 10,-pi/2)};
  \end{axis}
\end{tikzpicture}

%%% ----------------------------------------------------------- pages 18--26
%%% Untitled counterparts of pages 1, 2, 4, 5, 7, 8, 10, 11 and 12.
%%% page 18: sin, untitled
\begin{tikzpicture}
  \begin{axis}[sincos]
    \addplot[very thick] {sin(x)};
  \end{axis}
\end{tikzpicture}

%%% page 19: sin restricted, untitled
\begin{tikzpicture}
  \begin{axis}[
    sincos,
    xtick={-pi/2, pi/2},
    xticklabels={\(-\pi/2\), \(\pi/2\)},
    ]
    \addplot[very thick, main, domain={-pi/2}:{pi/2}] {sin(x)};
  \end{axis}
\end{tikzpicture}

%%% page 20: cos, untitled
\begin{tikzpicture}
  \begin{axis}[sincos]
    \addplot[very thick] {cos(x)};
  \end{axis}
\end{tikzpicture}

%%% page 21: cos restricted, untitled
\begin{tikzpicture}
  \begin{axis}[sincos]
    \addplot[very thick, main, domain={0}:{pi}] {cos(x)};
  \end{axis}
\end{tikzpicture}

%%% page 22: tan, untitled
\begin{tikzpicture}
  \begin{axis}[tansec]
    \addplot[very thick] {tan(x)};
    \addplot[dotted] coordinates {(-1*pi/2, 5)  (-1*pi/2, -5)};
    \addplot[dotted] coordinates {(-3*pi/2, 5)  (-3*pi/2, -5)};
    \addplot[dotted] coordinates {( 1*pi/2, 5)  ( 1*pi/2, -5)};
    \addplot[dotted] coordinates {( 3*pi/2, 5)  ( 3*pi/2, -5)};
  \end{axis}
\end{tikzpicture}

%%% page 23: tan restricted, untitled
\begin{tikzpicture}
  \begin{axis}[
    tansec,
    xtick={-pi/2, pi/2},
    xticklabels={\(-\pi/2\), \(\pi/2\)},
    ]
    \addplot[very thick, main, domain={-pi/2}:{pi/2}] {tan(x)};
    \addplot[dotted, supp] coordinates {(-1*pi/2, 5)  (-1*pi/2, -5)};
    \addplot[dotted, supp] coordinates {( 1*pi/2, 5)  ( 1*pi/2, -5)};
  \end{axis}
\end{tikzpicture}

%%% page 24: cot, untitled
\begin{tikzpicture}
  \begin{axis}[tansec]
    \addplot[very thick] {cot(x)};
    \addplot[dotted] coordinates {(-2*pi, 5)  (-2*pi, -5)};
    \addplot[dotted] coordinates {(-1*pi, 5)  (-1*pi, -5)};
    \addplot[dotted] coordinates {( 1*pi, 5)  ( 1*pi, -5)};
    \addplot[dotted] coordinates {( 2*pi, 5)  ( 2*pi, -5)};
  \end{axis}
\end{tikzpicture}

%%% page 25: cot restricted, untitled
\begin{tikzpicture}
  \begin{axis}[
    tansec,
    xtick={0,pi},
    xticklabels={\(0\), \(\pi\)}
    ]
    \addplot[very thick, main, domain={0}:{pi}] {cot(x)};
    \addplot[dotted, supp] coordinates {( 1*pi, 5)  ( 1*pi, -5)};
  \end{axis}
\end{tikzpicture}

%%% page 26: arccot, untitled
\begin{tikzpicture}
  \begin{axis}[
    arctan,
    ymin={0-pi/6}, ymax={pi+pi/6},
    ]
    \addplot[ultra thick, black] {pi/2-atan(x)};
    \addplot[dotted, supp] coordinates {(-10, 0)  ( 10, 0)};
    \addplot[dotted, supp] coordinates {(-10,pi)  ( 10,pi)};
  \end{axis}
\end{tikzpicture}

%%% ----------------------------------------------------------- pages 27--28
%%% Titled counterparts of pages 13 and 14.
%%% page 27: sec
\begin{tikzpicture}
  \begin{axis}[tansec, title={\(\sec(x)\)}]
    \addplot[very thick] {1/cos(x)};
    \addplot[dotted] coordinates {(-1*pi/2, 5)  (-1*pi/2, -5)};
    \addplot[dotted] coordinates {(-3*pi/2, 5)  (-3*pi/2, -5)};
    \addplot[dotted] coordinates {( 1*pi/2, 5)  ( 1*pi/2, -5)};
    \addplot[dotted] coordinates {( 3*pi/2, 5)  ( 3*pi/2, -5)};
  \end{axis}
\end{tikzpicture}

%%% page 28: csc
\begin{tikzpicture}
  \begin{axis}[tansec, title={\(\csc(x)\)}]
    \addplot[very thick] {1/sin(x)};
    \addplot[dotted] coordinates {(-1*pi, 5)  (-1*pi, -5)};
    \addplot[dotted] coordinates {(-2*pi, 5)  (-2*pi, -5)};
    \addplot[dotted] coordinates {( 1*pi, 5)  ( 1*pi, -5)};
    \addplot[dotted] coordinates {( 2*pi, 5)  ( 2*pi, -5)};
  \end{axis}
\end{tikzpicture}

%%% ----------------------------------------------------------- pages 29--31
%%% arcsin, arccos and arctan on one common y-range.  Same y-limits and
%%% "axis equal image" give the three pictures the same height, so placing
%%% them side by side puts every x-axis at the same height.
%%% page 29: arcsin, aligned
\begin{tikzpicture}
  \begin{axis}[
    arcsin,
    title={\(\theta = \arcsin(x)\) on domain \([-1,1]\)},
    ymin={-pi/2-pi/6}, ymax={pi+pi/6},
    ]
    \addplot[ultra thick, black] {asin(x)};
    \rangeseg{-pi/2}{pi/2}
    \rangedot{-pi/2}
    \rangedot{pi/2}
    \curveends{-1}{-pi/2}{1}{pi/2}
  \end{axis}
\end{tikzpicture}

%%% page 30: arccos, aligned
\begin{tikzpicture}
  \begin{axis}[
    arcsin,
    title={\(\theta = \arccos(x)\) on domain \([-1,1]\)},
    ymin={-pi/2-pi/6}, ymax={pi+pi/6},
    ]
    \addplot[ultra thick, black] {acos(x)};
    \rangeseg{0}{pi}
    \rangedot{0}
    \rangedot{pi}
    \curveends{-1}{pi}{1}{0}
  \end{axis}
\end{tikzpicture}

%%% page 31: arctan, aligned
\begin{tikzpicture}
  \begin{axis}[
    arctan,
    title={\(\theta = \arctan(x)\) on domain \((-\infty,\infty)\)},
    ymin={-pi/2-pi/6}, ymax={pi+pi/6},
    ]
    \addplot[ultra thick, black] {atan(x)};
    \addplot[dotted, supp] coordinates {(-10, pi/2)  ( 10, pi/2)};
    \addplot[dotted, supp] coordinates {(-10,-pi/2)  ( 10,-pi/2)};
    \rangeseg{-pi/2}{pi/2}
    \rangehole{-pi/2}
    \rangehole{pi/2}
  \end{axis}
\end{tikzpicture}
\end{document}
